% Proposta separada: preserva o fluxo e as contagens; estilo da Figura 2.1.
\begin{tikzpicture}[x=1cm,y=1cm,
  elo/.style={ntcaixa,minimum width=6.15cm,minimum height=1.35cm,inner sep=0pt},
  titulo/.style={anchor=west,align=left,text=corTinta,font=\fontsize{8.5}{10}\selectfont,inner sep=0pt,text width=4.85cm},
  detalhe/.style={anchor=west,align=left,text=corCinza,font=\fontsize{7.2}{9}\selectfont,inner sep=0pt,text width=4.85cm},
  numero/.style={circle,minimum size=5.1mm,inner sep=0pt,draw=corCinzaC,line width=.4pt,fill=white,text=corCinza,font=\fontsize{8}{9}\selectfont},
  saida/.style={anchor=west,align=left,text=corTinta,font=\fontsize{8.5}{10.5}\selectfont,inner sep=0pt,text width=5.35cm},
  motivo/.style={saida,text=corCinza,font=\fontsize{7.4}{9.5}\selectfont}]
  \begin{scope}[on background layer]
    \fill[corFundo,rounded corners=2pt] (.70,-9.03) rectangle (7.12,1.13);
    \draw[corCinzaC,line width=.5pt,rounded corners=2pt] (7.48,-9.03) rectangle (13.95,1.13);
  \end{scope}
  \node[titulo,font=\scriptsize\bfseries,text width=5.6cm] at (.99,.79) {Percurso de seleção};
  \node[titulo,font=\scriptsize\bfseries,text width=5.8cm] at (7.78,.79) {Remoções e exclusões};
  \draw[corCinzaC,line width=.35pt] (.99,.51) -- (6.82,.51);
  \draw[corCinzaC,line width=.35pt] (7.78,.51) -- (13.65,.51);

  \foreach \i/\y in {1/-.43,2/-2.08,3/-3.73,4/-5.38}
    \node[elo] (n\i) at (3.91,\y) {};
  \node[elo,draw=corFG,line width=.7pt,fill=corFG!6,minimum height=2.4cm] (n5) at (3.91,-7.62) {};
  \foreach \i in {1,2,3,4}
    \node[numero] at ([xshift=4.5mm]n\i.west) {\i};
  \node[numero,draw=corFG,text=corFG] at ([xshift=4.5mm,yshift=6.9mm]n5.west) {5};

  \node[titulo] at ([xshift=9mm,yshift=3.5mm]n1.west) {Identificação nas bases};
  \node[titulo,font=\fontsize{9}{10}\selectfont] at ([xshift=9mm,yshift=-.2mm]n1.west) {$n=1\,478$ registos};
  \node[detalhe] at ([xshift=9mm,yshift=-4.3mm]n1.west) {Scopus: 954 \enspace · \enspace Web of Science: 524};
  \node[titulo] at ([xshift=9mm,yshift=3mm]n2.west) {Triagem por título e resumo};
  \node[titulo,font=\fontsize{9}{10}\selectfont] at ([xshift=9mm,yshift=-2.5mm]n2.west) {$n=925$ registos};
  \node[titulo] at ([xshift=9mm,yshift=3mm]n3.west) {Relatórios procurados};
  \node[titulo,font=\fontsize{9}{10}\selectfont] at ([xshift=9mm,yshift=-2.5mm]n3.west) {para recuperação: $n=599$};
  \node[titulo] at ([xshift=9mm,yshift=3mm]n4.west) {Avaliação da elegibilidade};
  \node[titulo,font=\fontsize{9}{10}\selectfont] at ([xshift=9mm,yshift=-2.5mm]n4.west) {$n=459$ relatórios};
  \node[titulo,font=\fontsize{8.5}{10}\selectfont\bfseries] at ([xshift=9mm,yshift=7mm]n5.west) {Incluídos --- Corpus A};
  \node[titulo,text=corFG,font=\fontsize{10}{11}\selectfont] at ([xshift=9mm,yshift=2.4mm]n5.west) {$n=331$ publicações};
  \draw[corFG!25,line width=.35pt] ([xshift=9mm,yshift=-.8mm]n5.west) -- ([xshift=-3mm,yshift=-.8mm]n5.east);
  \node[detalhe] at ([xshift=9mm,yshift=-4.4mm]n5.west) {Corpus B: 40, das quais 39 cartografadas};
  \node[detalhe] at ([xshift=9mm,yshift=-8.6mm]n5.west) {Casos exploratórios: 13, fora de B};

  \foreach \a/\b in {1/2,2/3,3/4,4/5}
    \draw[ntseta] (n\a.south) -- (n\b.north);
  \foreach \i in {1,2,3,4}
    \draw[ntseta] (n\i.east) -- (7.80,0 |- n\i.center);
  \node[saida] at (7.96,-.22) {Duplicados removidos: $n=553$};
  \node[motivo] at (7.96,-.65) {Antes da triagem};
  \node[saida] at (7.96,-1.86) {Registos excluídos: $n=326$};
  \node[motivo] at (7.96,-2.29) {Fora de âmbito na triagem};
  \node[saida] at (7.96,-3.50) {Não recuperados: $n=137$};
  \node[motivo] at (7.96,-3.96) {Duplicados detetados tardiamente: $n=3$};
  \node[saida] at (7.96,-4.85) {Relatórios excluídos: $n=128$};
  \node[motivo] at (7.96,-5.30) {Sem modelo com detalhe suficiente: $n=117$};
  \node[motivo] at (7.96,-5.71) {Idioma não elegível: $n=10$};
  \node[motivo] at (7.96,-6.12) {Retratação: $n=1$};
  \draw[corCinzaC,line width=.35pt] (7.96,-6.42) -- (13.60,-6.42);
  \node[motivo,font=\fontsize{7.5}{10}\selectfont,text width=5.35cm] at (7.96,-7.03)
    {Rastreio de referências não executado.\\Previsto no protocolo (D9).};
  \node[motivo,font=\fontsize{7.5}{10}\selectfont,text width=5.35cm] at (7.96,-8.08)
    {B é um subconjunto de A.\\Os 13 casos exploratórios integram A,\\mas não as contagens de B.};
\end{tikzpicture}
